\documentclass[12pt, a4paper, twoside]{book}

% ─── Packages ───────────────────────────────────────────────────────────────
\usepackage[a4paper, inner=3.0cm, outer=2.5cm, top=2.8cm, bottom=2.8cm, includeheadfoot, headheight=24pt]{geometry}
\usepackage[T1]{fontenc}
\usepackage[utf8]{inputenc}
\usepackage{amsmath, amsthm}
\usepackage[charter, cal=cmcal]{mathdesign}
\usepackage[scaled=0.92]{helvet}
\usepackage{mathtools}
\usepackage{physics}
\usepackage{graphicx}
\usepackage{xcolor}
\usepackage{tikz}
\usetikzlibrary{arrows.meta, decorations.pathmorphing, shapes.geometric,
                positioning, calc, patterns, backgrounds, decorations.markings}
\usepackage{pgfplots}
\pgfplotsset{compat=1.18}
\usepackage{tcolorbox}
\tcbuselibrary{skins, breakable, theorems}
\usepackage{enumitem}
\usepackage{booktabs}
\usepackage{array}
\usepackage{multirow}
\usepackage{fancyhdr}
\usepackage{titlesec}
\usepackage{hyperref}
\usepackage{wrapfig}
\usepackage{float}
\usepackage{caption}
\usepackage{subcaption}
\usepackage{epigraph}
\usepackage{lettrine}
\usepackage{mdframed}
\usepackage{setspace}
\usepackage{parskip}
\usepackage{siunitx}

% ─── Colours ────────────────────────────────────────────────────────────────
\definecolor{deep}{HTML}{0F172A}      % Deep Slate/Charcoal for body text/titles
\definecolor{accent}{HTML}{E11D48}    % Crimson/Coral for problem counters
\definecolor{teal}{HTML}{0D9488}      % Electric Teal for rules and highlights
\definecolor{gold}{HTML}{D97706}      % Gold for puzzle borders
\definecolor{midblue}{HTML}{334155}   % Muted Blue-Gray for section headings
\definecolor{purple}{HTML}{6D28D9}    % Purple for reflection borders
\definecolor{lightbg}{HTML}{F8FAFC}   % Very light gray-blue for exploration bg
\definecolor{warmbg}{HTML}{FEF3C7}    % Warm gold-yellow for puzzle bg
\definecolor{greenbg}{HTML}{F0FDF4}   % Light green for thought/keyidea bg

% ─── Page geometry fix ──────────────────────────────────────────────────────
\setlength{\headheight}{24pt}

% ─── Hyperref ───────────────────────────────────────────────────────────────
\hypersetup{
  colorlinks=true, linkcolor=midblue, urlcolor=accent, citecolor=teal,
  pdftitle={Discovery Physics: Journey 2 — Motion and Its Causes},
  pdfauthor={Discovery Institute},
}

% ─── Page style ─────────────────────────────────────────────────────────────
\pagestyle{fancy}
\fancyhf{}
\fancyhead[LE]{\small\textcolor{deep}{\textit{Discovery Physics}}}
\fancyhead[RO]{\small\textcolor{deep}{\textit{\leftmark}}}
\fancyfoot[C]{\small\textcolor{deep}{\thepage}}
\renewcommand{\headrulewidth}{0.4pt}
\renewcommand{\headrule}{\hbox to\headwidth{\color{teal}\leaders\hrule height \headrulewidth\hfill}}

% ─── Chapter / Section titles ───────────────────────────────────────────────
\titleformat{\chapter}[display]
  {\normalfont\huge\bfseries\color{deep}}
  {\textcolor{accent}{\chaptertitlename\ \thechapter}}{12pt}{\Huge}
  [\vspace{4pt}\textcolor{teal}{\titlerule[1.5pt]}]

\titleformat{\section}
  {\normalfont\Large\bfseries\color{midblue}}
  {\textcolor{accent}{\thesection}}{1em}{}

\titleformat{\subsection}
  {\normalfont\large\bfseries\color{deep}}
  {\textcolor{teal}{\thesubsection}}{1em}{}

% ─── Custom boxes ───────────────────────────────────────────────────────────
\newtcolorbox{puzzlebox}[1][]{
  enhanced, breakable,
  colback=warmbg, colframe=gold,
  fonttitle=\bfseries\large\color{deep}, title={The Puzzle},
  attach boxed title to top left={yshift=-2mm, xshift=4mm},
  boxed title style={colback=gold, colframe=gold}, #1}

\newtcolorbox{explobox}[1][]{
  enhanced, breakable,
  colback=lightbg, colframe=midblue,
  fonttitle=\bfseries\color{white}, coltitle=white, title={Exploration},
  attach boxed title to top left={yshift=-2mm, xshift=4mm},
  boxed title style={colback=midblue, colframe=midblue}, #1}

\newtcolorbox{discoverybox}[1][]{
  enhanced, breakable,
  colback=white, colframe=accent,
  fonttitle=\bfseries\large\color{white}, coltitle=white, title={Discovery},
  attach boxed title to top left={yshift=-2mm, xshift=4mm},
  boxed title style={colback=accent, colframe=accent}, drop shadow, #1}

\newtcolorbox{namebox}[1][]{
  enhanced,
  colback=deep!5, colframe=deep,
  fonttitle=\bfseries\color{white}, coltitle=white, title={What This Is Called},
  attach boxed title to top left={yshift=-2mm, xshift=4mm},
  boxed title style={colback=deep, colframe=deep}, #1}

\newtcolorbox{reflectionbox}[1][]{
  enhanced, breakable,
  colback=purple!5, colframe=purple,
  fonttitle=\bfseries\color{white}, coltitle=white, title={Reflection}, #1}

\newtcolorbox{historybox}[1][]{
  enhanced, breakable,
  colback=gold!8, colframe=gold!70!brown,
  fonttitle=\bfseries\color{deep}, title={Historical Note}, #1}

% ─── Problem counter ────────────────────────────────────────────────────────
\newcounter{problem}[chapter]
\newcommand{\problem}[1][]{%
  \refstepcounter{problem}%
  \medskip\noindent
  \textcolor{accent}{\textbf{Problem \theproblem.}}\ %
  \ifx&#1&\else\textit{(#1)}\ \fi}

\newcommand{\keyidea}[1]{%
  \begin{tcolorbox}[enhanced, colback=teal!8, colframe=teal,
    left=4mm, fontupper=\bfseries\color{deep}]#1\end{tcolorbox}}

% ─── Theorem environments ───────────────────────────────────────────────────
\newtheoremstyle{disc}{8pt}{8pt}{\itshape}{}{\bfseries\color{midblue}}{.}{.5em}{}
\theoremstyle{disc}
\newtheorem{principle}{Principle}[chapter]
\newtheorem{question}{Question}[chapter]

% ─── Spacing ────────────────────────────────────────────────────────────────
\setlength{\epigraphwidth}{0.65\textwidth}
\renewcommand{\epigraphflush}{center}
\setlength{\parskip}{6pt}
\setlength{\parindent}{0pt}
\onehalfspacing

% ════════════════════════════════════════════════════════════════════════════
\begin{document}


\setcounter{chapter}{1}
\setcounter{section}{4}

\section{Module 2.5 — Conservation of Momentum}
\label{sec:momentum_conservation}

\epigraph{
  \itshape In physics, symmetry is more fundamental than substance.
}{--- Frank Wilczek}

\subsection{The Central Question}

\begin{tcolorbox}[colback=deep!5, colframe=deep, fontupper=\large\bfseries\color{deep}]
When two objects interact, what quantity is always the same before and after? And why should anything be conserved at all?
\end{tcolorbox}

\subsection{Discovering Momentum}

\begin{figure}[H]
  \centering
  \includegraphics[width=0.9\textwidth]{figures/fig_j2_momentum}
  \caption{Left: an elastic collision between two equal masses. The first mass stops; the second moves off with the first's velocity. Both momentum and kinetic energy are conserved. Right: momentum bars before and after confirm the total is unchanged.}
  \label{fig:momentum}
\end{figure}

Consider two objects colliding. We want to find a quantity that remains the same before and after the collision, regardless of the details of the collision.

\begin{explobox}
\problem Two carts on a frictionless track collide. Use the data below to find a conserved quantity.

\begin{center}
\begin{tabular}{lccccc}
\toprule
Experiment & $m_1$ & $v_{1i}$ & $m_2$ & $v_{2i}$ & $v_{1f}$, $v_{2f}$ \\
\midrule
1 & $\SI{1}{kg}$ & $\SI{4}{m/s}$ & $\SI{1}{kg}$ & $0$ & $0$, $\SI{4}{m/s}$ \\
2 & $\SI{2}{kg}$ & $\SI{3}{m/s}$ & $\SI{1}{kg}$ & $0$ & $\SI{1}{m/s}$, $\SI{4}{m/s}$ \\
3 & $\SI{1}{kg}$ & $\SI{5}{m/s}$ & $\SI{2}{kg}$ & $-\SI{1}{m/s}$ & $-\SI{1}{m/s}$, $\SI{3}{m/s}$ \\
\bottomrule
\end{tabular}
\end{center}
Compute $m_1 v_{1i} + m_2 v_{2i}$ and $m_1 v_{1f} + m_2 v_{2f}$ for each experiment. What do you find?

\problem Derive momentum conservation directly from Newton's third law. (Hint: during a collision, the force on object 1 from object 2 is equal and opposite to the force on object 2 from object 1. Apply $F = ma$ to both, then add the equations.)

\problem A rifle of mass $\SI{3}{kg}$ fires a bullet of mass $\SI{0.01}{kg}$ at $\SI{900}{m/s}$. What is the recoil velocity of the rifle?
\end{explobox}

\begin{namebox}
The conserved quantity is \textbf{momentum}: $\vec{p} = m\vec{v}$.

The total momentum of a system of objects is conserved---remains constant---whenever the objects interact only with each other (no external forces).

\textbf{Conservation of momentum} is one of the most universal laws in physics. It holds in quantum mechanics, in special relativity, in particle physics. It has never been observed to fail.
\end{namebox}

\subsection{Why Is Momentum Conserved? A Deeper Answer}

The derivation from Newton's third law shows momentum is conserved in Newtonian mechanics. But this is not the deepest answer.

In 1915, Emmy Noether proved one of the most beautiful theorems in all of physics:

\begin{tcolorbox}[colback=purple!5, colframe=purple, fonttitle=\bfseries\color{white},
  coltitle=white, title={Noether's Theorem (preview)}]
Every continuous symmetry of the laws of nature corresponds to a conserved quantity.

\medskip
Specifically:
\begin{itemize}
  \item \textbf{Translation symmetry in space} (physics is the same everywhere) $\implies$ conservation of \textbf{momentum}
  \item \textbf{Translation symmetry in time} (physics is the same always) $\implies$ conservation of \textbf{energy}
  \item \textbf{Rotation symmetry} (physics is the same in all directions) $\implies$ conservation of \textbf{angular momentum}
\end{itemize}
\end{tcolorbox}

Momentum is conserved not because of a coincidence in Newton's laws, but because \textit{space itself is the same in all locations}. If physics were different in Kenya than in Japan, momentum would not be conserved. The conservation law encodes a symmetry of the universe.

You will encounter Noether's theorem properly in Mathematics Journey 4 Module 5 and Physics Journey 6 Module 4. For now, simply hold this thought: conservation laws are not accidents. They are shadows of symmetry.

\begin{reflectionbox}
\textbf{Reflection.} You have now encountered two conservation laws: conservation of momentum and (earlier) conservation of energy. Both emerged from specific problems. But Noether's theorem says they are both shadows of symmetry.

Write a short paragraph: what does it mean for the universe to be ``symmetric'' in some direction? Can you think of situations where this symmetry might break down---and what would that imply for conservation laws?
\end{reflectionbox}

% ────────────────────────────────────────────────────────────────────────────

\end{document}
